\label{anlage:katalog}

Diese Anlage dokumentiert die im Berichtszeitraum entwickelten Oberflächen- und Navigationsstände des Bauteilkatalogs (Abschnitt~1.2.3.2). Die Screenshots zeigen Zwischenstände der Benutzeroberfläche und keinen abgeschlossenen Funktionsumfang.

\begin{Figure}[title=Bauteilobjekt, height=17cm, image-fit=contain, anchor=center]
\SemioImage{asset/katalog/bauteilobjekt.png}
\end{Figure}

Die Katalognavigation führt von der Typologienübersicht über die Typen einer Struktur bis zu den einzelnen Bauteilen (Pieces). Das aktive Bauteilobjekt ABAUF-CB-001 zeigt Anschlüsse, den Nachweisstatus je Fachbereich (Geometrie, Statik, Brandschutz, System, CO\textsubscript{2}) sowie den Reifegrad.

\begin{Figure}[title=Filteransicht, height=9cm, image-fit=contain, anchor=center]
\SemioImage{asset/katalog/filteransicht.png}
\end{Figure}

Das Filter-Panel gliedert Filterkriterien nach Typologie, Materialfamilie, Geometrie, Funktion im Entwurf, Verfügbarkeit, Reifegrad sowie Risiko und Nachweise. Aktive Filter werden als Chips oberhalb der Kriteriengruppen zusammengefasst.

\begin{Figure}[title=Bauteilkarte, height=17cm, image-fit=contain, anchor=center]
\SemioImage{asset/katalog/bauteilkarte.png}
\end{Figure}

Der hierarchische Katalogbaum gliedert Bauteile nach Systemebene (Structure, Envelope, Interior) und darunterliegenden Typen; die Bauteilkarte am unteren Rand bündelt Geometrie, Anschlüsse und Nachweise des ausgewählten Bauteils.
